% from ega0-22-fr.tex
\oldpage[0\textsubscript{IV-1}]{187}

rang fini~; nous considérons les triplets $(A,\alpha,\beta)$, où $A$ est une
$k$-algèbre locale noethérienne complète \emph{formellement lisse} d'idéal
maximal $\mathfrak m$, $\alpha:A/\mathfrak m\xrightarrow{\sim}K$ un
isomorphisme de $k$-algèbres,
$\beta:\mathfrak m/\mathfrak m^2\xrightarrow{\sim}V$ un di-isomorphisme
d'espaces vectoriels (pour l'isomorphisme $\alpha$). Étant donnés deux tels
triplets $(A,\alpha,\beta)$, $(A',\alpha',\beta')$, une \emph{équivalence} du
premier sur le second est un $k$-isomorphisme $u:A\xrightarrow{\sim}A'$ tel
que (si $\mathfrak m'$ est l'idéal maximal de $A'$), les diagrammes
\[
\xymatrix@C=4em@R=3em{
  A/\mathfrak m \ar[r]^{\operatorname{gr}^0(u)} \ar[d]_{\alpha} &
  A'/\mathfrak m' \ar[d]^{\alpha'} \\
  K \ar[r]_{1_K} & K
}
\qquad
\xymatrix@C=4em@R=3em{
  \mathfrak m/\mathfrak m^2 \ar[r]^{\operatorname{gr}^1(u)}
    \ar[d]_{\beta} &
  \mathfrak m'/\mathfrak m'^2 \ar[d]^{\beta'} \\
  V \ar[r]_{1_V} & V
}
\]
sont commutatifs. Il est immédiat que les \emph{classes d'équivalence} des
triplets $(A,\alpha,\beta)$ forment un \emph{ensemble} que nous noterons
$\operatorname{FL}(K/k,V)$. Remarquons maintenant qu'à tout triplet
$(A,\alpha,\beta)$ est associé le $K$-homomorphisme composé d'espaces
vectoriels
\[
  \Upsilon_{K/k}\xrightarrow{\chi_{A/k}}
  \mathfrak m/\mathfrak m^2\xrightarrow{\beta}V
\]
(où $\chi_{A/k}$ est considéré comme un di-homomorphisme relatif à
$\alpha^{-1}$)~; la définition précédente prouve (par (20.6.22.2)) que ce
$K$-homomorphisme \emph{ne dépend que de la classe d'équivalence} de
$(A,\alpha,\beta)$ dans $\operatorname{FL}(K/k,V)$~; autrement dit on a
défini une application canonique
\begin{equation}
\tag{22.2.4.1}
\label{0.22.2.4.1-fr}
  \operatorname{FL}(K/k,V)\longrightarrow
  \operatorname{Hom}_K(\Upsilon_{K/k},V).
\end{equation}
\end{env}

\begin{proposition}[22.2.5]
\label{0.22.2.5-fr}
Soient $k$ un corps, $K$ une extension de $k$, $V$ un $K$-espace vectoriel
de rang fini. L'application canonique (22.2.4.1) est une bijection de
$\operatorname{FL}(K/k,V)$ sur l'ensemble des $K$-homomorphismes injectifs
de $\Upsilon_{K/k}$ dans $V$.
\end{proposition}

\begin{proof}
(i) Pour montrer que (22.2.4.1) est injective, il faut prouver ce qui suit~:
soient $A$, $A'$ deux $k$-algèbres locales noethériennes complètes
formellement lisses, $\mathfrak m$, $\mathfrak m'$ leurs idéaux maximaux
respectifs, $\alpha:A/\mathfrak m\xrightarrow{\sim}K$,
$\alpha':A'/\mathfrak m'\xrightarrow{\sim}K$ deux $k$-isomorphismes~;
supposons donnés deux $k$-isomorphismes
\[
  v^0:A/\mathfrak m\xrightarrow{\sim}A'/\mathfrak m',\qquad
  v^1:\mathfrak m/\mathfrak m^2\xrightarrow{\sim}
      \mathfrak m'/\mathfrak m'^2
\]
tels que les diagrammes
\begin{equation}
\tag{22.2.5.1}
\label{0.22.2.5.1-fr}
\xymatrix@C=4em@R=3em{
  A/\mathfrak m \ar[r]^{v^0} \ar[dr]_{\alpha} &
  A'/\mathfrak m' \ar[dl]^{\alpha'} \\
  & K
}
\qquad
\xymatrix@C=4em@R=3em{
  \mathfrak m/\mathfrak m^2 \ar[r]^{v^1} &
  \mathfrak m'/\mathfrak m'^2 \\
  \Upsilon_{K/k} \ar[u]^{\chi_A} \ar[ur]_{\chi_{A'}} &
}
\end{equation}
soient commutatifs. Il s'agit de prouver qu'il existe un $k$-isomorphisme
$u:A\xrightarrow{\sim}A'$ tel que $\operatorname{gr}^0(u)=v^0$ et
$\operatorname{gr}^1(u)=v^1$. Remarquons en premier lieu que puisque $K$ est
un corps, l'homomorphisme canonique (20.6.7.1) est \emph{bijectif} (20.6.11),
donc il existe déjà un $k$-isomorphisme
$v:A/\mathfrak m^2\xrightarrow{\sim}A'/\mathfrak m'^2$ tel que
$\operatorname{gr}^0(v)=v^0$ et $\operatorname{gr}^1(v)=v^1$. Notons
maintenant que puisque $A'$ est métrisable et complète, et $A$ une
$k$-algèbre formellement lisse, l'homomorphisme composé
$A\to A/\mathfrak m^2\xrightarrow{v}A'/\mathfrak m'^2$ se factorise en
$A\xrightarrow{u}A'\to A'/\mathfrak m'^2$, où $u$ est

